% s4_results.tex — Results section
% Paper: "Tokenized Housing and Lifecycle Portfolio Choice:
%         A Decoupling of Location from Housing Exposure"
%
% Source: docs/welfare_decomp_v4.md (CEV formula, table shells, channel plan)
%         docs/sensitivity_grid_v4.md (sweep design)
%         paper/outline_v4.md (section 4 structure)
%
% Status: SKELETON — numerical placeholders marked [PLACEHOLDER] throughout.
%         All table shells match the pre-registered spec in welfare_decomp_v4.md.
%         Figures 1 and 2 are placeholder stubs.
%         Fill numerical entries once output/diagnostics/p6_option1_*.json land.
%
% Compilation: included from main.tex via % s4_results.tex — Results section
% Paper: "Tokenized Housing and Lifecycle Portfolio Choice:
%         A Decoupling of Location from Housing Exposure"
%
% Source: docs/welfare_decomp_v4.md (CEV formula, table shells, channel plan)
%         docs/sensitivity_grid_v4.md (sweep design)
%         paper/outline_v4.md (section 4 structure)
%
% Status: SKELETON — numerical placeholders marked [PLACEHOLDER] throughout.
%         All table shells match the pre-registered spec in welfare_decomp_v4.md.
%         Figures 1 and 2 are placeholder stubs.
%         Fill numerical entries once output/diagnostics/p6_option1_*.json land.
%
% Compilation: included from main.tex via % s4_results.tex — Results section
% Paper: "Tokenized Housing and Lifecycle Portfolio Choice:
%         A Decoupling of Location from Housing Exposure"
%
% Source: docs/welfare_decomp_v4.md (CEV formula, table shells, channel plan)
%         docs/sensitivity_grid_v4.md (sweep design)
%         paper/outline_v4.md (section 4 structure)
%
% Status: SKELETON — numerical placeholders marked [PLACEHOLDER] throughout.
%         All table shells match the pre-registered spec in welfare_decomp_v4.md.
%         Figures 1 and 2 are placeholder stubs.
%         Fill numerical entries once output/diagnostics/p6_option1_*.json land.
%
% Compilation: included from main.tex via % s4_results.tex — Results section
% Paper: "Tokenized Housing and Lifecycle Portfolio Choice:
%         A Decoupling of Location from Housing Exposure"
%
% Source: docs/welfare_decomp_v4.md (CEV formula, table shells, channel plan)
%         docs/sensitivity_grid_v4.md (sweep design)
%         paper/outline_v4.md (section 4 structure)
%
% Status: SKELETON — numerical placeholders marked [PLACEHOLDER] throughout.
%         All table shells match the pre-registered spec in welfare_decomp_v4.md.
%         Figures 1 and 2 are placeholder stubs.
%         Fill numerical entries once output/diagnostics/p6_option1_*.json land.
%
% Compilation: included from main.tex via \input{sections/s4_results}
%              Requires amsmath, booktabs, multirow, siunitx, graphicx, caption.

% ─────────────────────────────────────────────────────────────────────────────
\section{Results}
\label{sec:results}
% ─────────────────────────────────────────────────────────────────────────────

We report the baseline welfare results in Section~\ref{sec:results:baseline},
the channel decomposition in Section~\ref{sec:results:decomp}, the
pre-registered falsification tests in Section~\ref{sec:results:falsification},
and sensitivity analysis in Section~\ref{sec:results:sensitivity}.

All welfare gains are expressed as consumption equivalent variations (CEV) at the
initial state
$\mathbf{s}_0 = (t{=}1,\,w_{\text{mid}},\,z_{\text{mid}},\,\ell{=}A,\,x_{A,0}{=}0,\,x_{B,0}{=}0)$:
\begin{equation}
    \text{CEV}(R_a \text{ vs } R_b) \;=\;
    \left(\frac{V_1^{R_a}(\mathbf{s}_0)}{V_1^{R_b}(\mathbf{s}_0)}\right)^{1/(1-\gamma)}
    - 1,
    \label{eq:cev}
\end{equation}
where $V_1^{R}(\mathbf{s}_0)$ is the time-1 value function under regime $R$
evaluated at $\mathbf{s}_0$.%
\footnote{Both value functions are negative under CRRA with finite wealth,
so the ratio is positive and the formula is well-defined for $\gamma \neq 1$.
The CEV measures the fraction of lifetime consumption in regime $R_b$ that the
household would give up to gain permanent access to regime $R_a$.}
Appendix~\ref{app:numerical} reports convergence diagnostics and full-grid
robustness.

% ─────────────────────────────────────────────────────────────────────────────
\subsection{Baseline Welfare Results}
\label{sec:results:baseline}
% ─────────────────────────────────────────────────────────────────────────────

The headline welfare result is reported in Table~\ref{tab:cev_baseline}.
The first row gives $\text{CEV}(\text{E2\_2L vs E1\_2L})$, the lifetime
welfare value of tokenization relative to traditional binary homeownership.
Rows 2--4 decompose this gain into three channels (detailed in
Section~\ref{sec:results:decomp}).
Row 5 reports $\text{CEV}(\text{E0 vs E1\_2L})$ as a sanity-check benchmark:
the welfare cost of pure renting relative to traditional ownership.

\begin{table}[t]
\centering
\caption{Welfare Value of Tokenization: Baseline and Robustness Panel
         (CEV, percentage points of lifetime consumption)}
\label{tab:cev_baseline}
\small
\begin{tabular}{lcccccc}
\toprule
 & Baseline & $\rho_{AB}=0$ & $\rho_{AB}=0.50$ & $\rho_{AB}=0.95$
 & $p_{\text{reloc}}=0$ & $p_{\text{reloc}}=0.12$ \\
\midrule
$\text{CEV}(\text{E2\_2L vs E1\_2L})$
    & \textsc{[P]} & \textsc{[P]} & \textsc{[P]} & \textsc{[P]}
    & \textsc{[P]} & \textsc{[P]} \\[2pt]
\quad Avoided-tx channel
    & \textsc{[P]} & \textsc{[P]} & \textsc{[P]} & \textsc{[P]}
    & \textsc{[P]} & \textsc{[P]} \\
\quad Maintained-hedge channel
    & \textsc{[P]} & \textsc{[P]} & \textsc{[P]} & \textsc{[P]}
    & \textsc{[P]} & \textsc{[P]} \\
\quad Cross-term $\xi$
    & \textsc{[P]} & \textsc{[P]} & \textsc{[P]} & \textsc{[P]}
    & \textsc{[P]} & \textsc{[P]} \\[4pt]
$\text{CEV}(\text{E0 vs E1\_2L})$ (renter benchmark)
    & \textsc{[P]} & --- & --- & --- & --- & --- \\
\addlinespace
Mean $\bar{x}_B$ at $\ell{=}A$
    & \textsc{[P]} & \textsc{[P]} & \textsc{[P]} & \textsc{[P]}
    & \textsc{[P]} & \textsc{[P]} \\
\bottomrule
\end{tabular}
\vspace{4pt}
\begin{minipage}{\linewidth}
\small\textit{Notes:}
Baseline: $\gamma=5$, $\beta=0.96$, $\rho_{AB}=0.50$,
$p_{\text{reloc,work}}=0.06$, $\tau_{\text{sell}}=0.06$,
$\tau_{\text{buy}}=0.025$, $\tau_{\text{token}}=0.01$.
Grid: $N_W=15$, $N_Z=5$, $N_{x_{\text{prev}}}=3$ (coarse; full-grid
robustness at $N_W=40$, $N_Z=9$, $N_{x_{\text{prev}}}=5$ in
Appendix~\ref{app:numerical}).
\textsc{[P]} = placeholder pending server1 baseline runs.
\end{minipage}
\end{table}

The primary diagnostic for mechanism activation is $\bar{x}_B$ at $\ell{=}A$:
the mean portfolio share of location-$B$ tokens held by households residing
at location $A$.
Under E1\_2L, $\bar{x}_B = 0$ by admissibility (binary tenure forbids
non-occupancy-location ownership).
Under E2\_2L, a positive $\bar{x}_B > 0$ at $\ell{=}A$ indicates that
households voluntarily pre-accumulate location-$B$ tokens before relocation ---
the mechanism that is uniquely enabled by token portability and the 6D state
extension (Option~1).
This is the key discriminator from the v3 Option~3 approximation, under which
$\bar{x}_B = 0$ held at all mobility rates.

Figure~\ref{fig:lifecycle_profiles} illustrates the mechanism at the
lifecycle level: the age profiles of mean $x_A$ and $x_B$ holdings
under E1\_2L and E2\_2L.
The key prediction is that $\bar{x}_B$ rises during working years (ages 25--65),
when $p_{\text{reloc,work}} = 0.06$, and declines post-retirement when the
relocation probability falls to $p_{\text{reloc,ret}} = 0.02$.
This pattern --- confirmed or denied by the server1 policy-array output ---
constitutes the strongest visual evidence for the pre-buying hedge mechanism.

\begin{figure}[t]
\centering
%% PLACEHOLDER — replace with lifecycle profile plots once server1 policy arrays land.
%% See paper/exhibit_memos/fig1_lifecycle_profiles.md for full production spec.
\fbox{%
  \begin{minipage}{0.88\linewidth}
  \vspace{2.5cm}
  \centering
  {\sffamily\small [PLACEHOLDER: Figure~1 — Lifecycle token-holding profiles.
  Panel A (E1\_2L): mean $x_A$ (solid) and $x_B = 0$ (dashed) over ages 25--80.
  Panel B (E2\_2L): mean $x_A$ (solid) and mean $x_B > 0$ (dashed) over ages 25--80.
  Source: per-age mean of \texttt{xA\_policy} and \texttt{xB\_policy} arrays at $\ell{=}A$,
  averaged over feasible $(w, z, x_{\text{prev}})$ states.
  Script: \texttt{scripts/plot\_lifecycle\_profiles.py}.]}
  \vspace{2.5cm}
  \end{minipage}%
}
\caption{Lifecycle mean token-holding profiles under traditional ownership (E1\_2L,
         left panel) and tokenized ownership (E2\_2L, right panel) at location $A$.
         E2\_2L households pre-accumulate location-$B$ tokens during working years
         (positive $\bar{x}_B$), saving the buying cost $\tau_{\text{buy}}$ on
         anticipated future relocation.
         E1\_2L households hold $\bar{x}_B = 0$ at all ages by admissibility.
         \textsc{[P]} = placeholder pending server1 run.}
\label{fig:lifecycle_profiles}
\end{figure}

% ─────────────────────────────────────────────────────────────────────────────
\subsection{Channel Decomposition}
\label{sec:results:decomp}
% ─────────────────────────────────────────────────────────────────────────────

We decompose the total welfare gain using three regime runs.
Let E1\_2L$_{\text{NT}}$ denote the no-transaction-cost counterfactual
($\tau_{\text{sell}} = \tau_{\text{buy}} = 0$, binary tenure otherwise).
The decomposition is:
\begin{align}
    \underbrace{\text{CEV}(\text{E2\_2L vs E1\_2L})}_{\text{total}}
    &\;=\;
    \underbrace{\text{CEV}(\text{E1\_2L}_{\text{NT}} \text{ vs E1\_2L})}_{\text{avoided-tx}}
    \;+\;
    \underbrace{\text{CEV}(\text{E2\_2L vs E1\_2L}_{\text{NT}})}_{\text{maintained-hedge}}
    \;+\; \xi,
    \label{eq:decomp}
\end{align}
where $\xi$ is the cross-term.
This decomposition is pre-registered in \texttt{docs/welfare\_decomp\_v4.md}
and is exact up to the cross-term, which was empirically near zero in v3
($\xi = 0.021$\%) and is expected to remain small in v4.

\begin{table}[t]
\centering
\caption{Channel Decomposition of the Tokenization Welfare Gain
         (CEV, percentage points of lifetime consumption)}
\label{tab:decomp}
\small
\begin{tabular}{lccl}
\toprule
Channel & CEV (\%) & Share of total & Description \\
\midrule
Total: $\text{CEV}(\text{E2\_2L vs E1\_2L})$
    & \textsc{[P]} & 100\% & Headline tokenization gain \\
\addlinespace
\quad (i)~Avoided-tx
    & \textsc{[P]} & \textsc{[P]}\%
    & $\tau_{\text{sell}} + \tau_{\text{buy}}$ burden of E1\_2L moves \\
\qquad of which $\tau_{\text{sell}}$
    & \textsc{[P]} & \textsc{[P]}\%
    & NAR selling commission (6\%) \\
\qquad of which $\tau_{\text{buy}}$
    & \textsc{[P]} & \textsc{[P]}\%
    & Buying closing costs (2.5\%) \\[4pt]
\quad (ii)~Maintained-hedge
    & \textsc{[P]} & \textsc{[P]}\%
    & Continuous-x + pre-buying hedge vs E1\_2L$_\text{NT}$ \\[4pt]
\quad (iii)~Cross-term $\xi$
    & \textsc{[P]} & \textsc{[P]}\%
    & Interaction; expected $\approx 0$ \\
\midrule
v3 Option~3 (no state ext.)
    & $+4.255$\% & (reference) & Previous fire's upper-bound approx \\
\quad of which avoided-tx
    & $+0.816$\% & 19.2\% & From v3 full-grid run \\
\quad of which continuous-x
    & $+3.411$\% & 80.2\% & Liu (2021) territory \\
\quad of which pre-buying hedge
    & $\approx 0$\% & $\approx 0$\% & Absent without state ext. (Option~3 limit) \\
\bottomrule
\end{tabular}
\vspace{4pt}
\begin{minipage}{\linewidth}
\small\textit{Notes:}
Upper panel: v4 Option~1 results (\textsc{[P]} = placeholder).
Lower panel: v3 Option~3 reference numbers (confirmed from server1 run,
commit 186da13).
The maintained-hedge channel in v4 encompasses both the continuous-x
subchannel (present in v3) and the new pre-buying hedge subchannel
enabled by the 6D state; a fourth regime decomposition
($\text{E2\_2L}_{\text{NT}}$) can further separate the two if H1 is confirmed.
\end{minipage}
\end{table}

The maintained-hedge channel in v4 captures two sub-mechanisms:
\begin{enumerate}
    \item \emph{Continuous-$x$ sub-channel} (present in v3): continuous
    fractional ownership of the occupied unit avoids the binary kink at
    $x_\ell = 1$, delivering a Liu~(2021)-style indivisibility relaxation.
    \item \emph{Pre-buying hedge sub-channel} (new in v4): by tracking
    $(x_{A,t-1}, x_{B,t-1})$ in the state, households face a per-period
    buying cost $\tau_{\text{buy}} \cdot \max(\Delta x, 0)$ and can
    pre-accumulate location-$B$ tokens at location $A$, arriving at $B$
    with $x_{B,t-1} > 0$ and avoiding the lump-sum purchase cost.
    The expected per-period premium per unit held is
    $p_{\text{reloc}} \times \tau_{\text{buy}} = 0.06 \times 0.025 = 0.15$\%
    (see Appendix~\ref{app:proof}).
\end{enumerate}

% ─────────────────────────────────────────────────────────────────────────────
\subsection{Falsification Tests}
\label{sec:results:falsification}
% ─────────────────────────────────────────────────────────────────────────────

We conduct three pre-registered falsification tests designed to identify
the pre-buying hedge channel.
These tests were pre-registered in \texttt{docs/welfare\_decomp\_v4.md}
before server1 runs were executed.
Crucially, the analogous tests \emph{failed} under v3 Option~3 (which lacked
the state extension): $\bar{x}_B$ remained zero at all mobility rates,
and the CEV was insensitive to $p_{\text{reloc}}$ and $\rho_{AB}$.

\begin{table}[t]
\centering
\caption{Pre-Registered Falsification Tests
         (Deviations from Baseline)}
\label{tab:falsification}
\small
\begin{tabular}{llcccc}
\toprule
Test & Counterfactual & $\text{CEV}$ & $\bar{x}_B$ at $\ell{=}A$
     & Pass criterion & Result \\
\midrule
(r) Relocation disabled
    & $p_{\text{reloc}} = 0$
    & \textsc{[P]}
    & \textsc{[P]}
    & $\bar{x}_B \to 0$;\; $\Delta\text{CEV} < -0.5$pp
    & \textsc{[P]} \\[4pt]
(m) Perfect correlation
    & $\rho_{AB} = 0.95$
    & \textsc{[P]}
    & \textsc{[P]}
    & $\bar{x}_B$ decreases;\; $\Delta\text{CEV} < 0$
    & \textsc{[P]} \\[4pt]
(q) No buying cost
    & $\tau_{\text{buy}} = 0$
    & \textsc{[P]}
    & \textsc{[P]}
    & Hedge channel $\to 0$;\; $\bar{x}_B \to 0$
    & \textsc{[P]} \\
\bottomrule
\end{tabular}
\vspace{4pt}
\begin{minipage}{\linewidth}
\small\textit{Notes:}
All tests compare to the baseline E2\_2L v4 run.
Pass criterion stated prior to server1 runs.
``Hedge channel'' in test (q) is
$\text{CEV}(\text{E2\_2L}_{\tau=0} \text{ vs E1\_2L}_{\text{NT}}) - \text{CEV}(\text{E2\_2L}_{\text{baseline}} \text{ vs E1\_2L}_{\text{NT}})$.
\textsc{[P]} = placeholder.
\end{minipage}
\end{table}

Test (r) shuts off the relocation shock entirely.
Without future relocation, there is no value to pre-accumulating
location-$B$ tokens while at $A$; the pre-buying hedge channel
should collapse.
If test (r) passes --- i.e., $\bar{x}_B \to 0$ when $p_{\text{reloc}} = 0$
--- it provides strong evidence that the positive $\bar{x}_B$ in the
baseline is driven by the pre-buying motive and not by a model artefact.

Test (m) raises the cross-location correlation to $\rho_{AB} = 0.95$,
nearly eliminating diversification benefit.
When location-$A$ and location-$B$ price processes are nearly identical,
holding $x_B$ at $\ell = A$ provides little hedge against diverging location returns.
The CEV and $\bar{x}_B$ should both decrease relative to the baseline.

Test (q) eliminates the buying cost ($\tau_{\text{buy}} = 0$), which
removes the financial incentive to pre-accumulate $x_B$ (there is no lump-sum
arrival cost to avoid).
The pre-buying hedge channel should collapse to zero.

% ─────────────────────────────────────────────────────────────────────────────
\subsection{Sensitivity Analysis}
\label{sec:results:sensitivity}
% ─────────────────────────────────────────────────────────────────────────────

Figure~\ref{fig:sensitivity_heatmap} depicts the total welfare gain
$\text{CEV}(\text{E2\_2L vs E1\_2L})$ as a function of the two primary
mechanism parameters: the cross-location correlation $\rho_{AB}$ and
the relocation probability $p_{\text{reloc}}$.
The full sweep design follows the pre-registered plan in
\texttt{docs/sensitivity\_grid\_v4.md}; only a subset is shown here.

\begin{figure}[t]
\centering
%% PLACEHOLDER — replace with actual heatmap once sweep results available
\fbox{%
  \begin{minipage}{0.72\linewidth}
  \vspace{2.5cm}
  \centering
  {\sffamily\small [PLACEHOLDER: Figure~2 — CEV sensitivity heatmap.
  X-axis: $\rho_{AB} \in \{0, 0.25, 0.50, 0.75, 0.95\}$.
  Y-axis: $p_{\text{reloc}} \in \{0, 0.02, 0.06, 0.12\}$.
  Heat colour: total CEV(\%), from dark (low) to bright (high).
  Source: sweep outputs from \texttt{scripts/sweep\_rhoAB.sh} and
  \texttt{scripts/sweep\_prelocate.sh} once server1 runs complete.]}
  \vspace{2.5cm}
  \end{minipage}%
}
\caption{Welfare sensitivity to cross-location correlation and relocation rate.
         Each cell reports $\text{CEV}(\text{E2\_2L vs E1\_2L})$ at
         $(\rho_{AB},\, p_{\text{reloc}})$.
         The mechanism prediction is a positive gradient in $p_{\text{reloc}}$
         (higher mobility $\Rightarrow$ higher hedge value) and a negative
         gradient in $\rho_{AB}$ (more correlated locations $\Rightarrow$
         smaller diversification benefit).}
\label{fig:sensitivity_heatmap}
\end{figure}

Table~\ref{tab:sensitivity} reports selected cross-sections of the
sensitivity grid for other key parameters.

\begin{table}[t]
\centering
\caption{Sensitivity of Headline CEV to Key Parameters
         (Total CEV, \%; baseline value on diagonal)}
\label{tab:sensitivity}
\small
\begin{tabular}{lcccccc}
\toprule
Parameter & & \multicolumn{5}{c}{Values} \\
\midrule
Risk aversion $\gamma$
    & & 3 & \textbf{5} & 8 & & \\
\quad $\text{CEV}(\text{E2\_2L vs E1\_2L})$
    & & \textsc{[P]} & \textsc{[P]} & \textsc{[P]} & & \\
\addlinespace
Buying cost $\tau_{\text{buy}}$
    & & 0\% & 1\% & \textbf{2.5\%} & 4\% & 6\% \\
\quad $\text{CEV}(\text{E2\_2L vs E1\_2L})$
    & & \textsc{[P]} & \textsc{[P]} & \textsc{[P]} & \textsc{[P]} & \textsc{[P]} \\
\addlinespace
Relocation rate $p_{\text{reloc,work}}$
    & & 0 & 0.02 & \textbf{0.06} & 0.12 & \\
\quad $\text{CEV}(\text{E2\_2L vs E1\_2L})$
    & & \textsc{[P]} & \textsc{[P]} & \textsc{[P]} & \textsc{[P]} & \\
\quad $\bar{x}_B$ at $\ell{=}A$
    & & \textsc{[P]} & \textsc{[P]} & \textsc{[P]} & \textsc{[P]} & \\
\addlinespace
Cross-location corr.\ $\rho_{AB}$
    & & 0 & 0.25 & \textbf{0.50} & 0.75 & 0.95 \\
\quad $\text{CEV}(\text{E2\_2L vs E1\_2L})$
    & & \textsc{[P]} & \textsc{[P]} & \textsc{[P]} & \textsc{[P]} & \textsc{[P]} \\
\addlinespace
\multicolumn{7}{l}{\textit{Asymmetric robustness (solver v4 extension, fire 22)}} \\
Location-$B$ return premium $\Delta\mu_B$
    & & $-1\%$ & \textbf{0\%} & $+1\%$ & & \\
\quad $\text{CEV}(\text{E2\_2L vs E1\_2L})$
    & & \textsc{[P]} & \textsc{[P]} & \textsc{[P]} & & \\
\quad $\bar{x}_B$ at $\ell{=}A$
    & & \textsc{[P]} & \textsc{[P]} & \textsc{[P]} & & \\[2pt]
Directional mobility ($p_{A\to B}$, $p_{B\to A}$)
    & & (0.06, 0.06) & (0.10, 0.03) & (0.10, 0.10) & & \\
\quad $\text{CEV}(\text{E2\_2L vs E1\_2L})$
    & & \textsc{[P]} & \textsc{[P]} & \textsc{[P]} & & \\
\addlinespace
\multicolumn{7}{l}{\textit{Mortgage activation (LTV constraint)}} \\
LTV ceiling $\theta_{\max}$
    & & \textbf{0} & 0.50 & 0.80 & & \\
\quad $\text{CEV}(\text{E2\_2L vs E1\_2L})$
    & & \textsc{[P]} & \textsc{[P]} & \textsc{[P]} & & \\
\bottomrule
\end{tabular}
\vspace{4pt}
\begin{minipage}{\linewidth}
\small\textit{Notes:}
Bold entries mark the baseline value.
One parameter varies at a time; all others held at baseline.
Sweep scripts: \texttt{scripts/sweep\_rhoAB.sh},
\texttt{scripts/sweep\_prelocate.sh}, \texttt{scripts/sweep\_txcost.sh},
\texttt{scripts/sweep\_asymmetric.sh} (location-$B$ return premium and directional mobility),
\texttt{scripts/sweep\_mortgage.sh} (LTV).
$\Delta\mu_B = \mu_{h,B} - \mu_{h,A}$: the excess Jensen-corrected log-mean return of
location $B$ relative to $A$.
Positive $\Delta\mu_B$ strengthens the pre-buying motive for $x_B$ by adding a return
premium to the tau\_buy-saving incentive.
\textsc{[P]} = placeholder pending server1 sweep runs.
\end{minipage}
\end{table}

Two monotonicity predictions follow from the model mechanism.
First, for $\tau_{\text{buy}}$: as the buying cost rises, the pre-accumulation
motive strengthens and $\bar{x}_B$ should increase, lifting the hedge sub-channel.
Second, for $p_{\text{reloc}}$: higher mobility raises the expected value of
pre-accumulated $x_B$ holdings; at $p_{\text{reloc}} = 0$ the hedge channel
collapses (falsification test r).
If these monotonicity patterns hold in the sweep results, they constitute
additional evidence for the pre-buying hedge mechanism.

% ─────────────────────────────────────────────────────────────────────────────
\subsection{Summary of Mechanism Evidence}
\label{sec:results:summary}
% ─────────────────────────────────────────────────────────────────────────────

The results establish \textsc{[P]} as the headline lifetime welfare gain from
tokenization under the baseline calibration.
The gain decomposes additively into an avoided-transaction-cost channel
(\textsc{[P]}\%) and a maintained-hedge channel (\textsc{[P]}\%), with
a near-zero cross-term.
The maintained-hedge channel is itself divisible into the continuous-$x$
sub-channel present in prior work (Liu, 2021) and the novel pre-buying hedge
sub-channel unlocked by the 6D state extension.

The three falsification tests provide mechanism identification.
Under all three counterfactuals --- zero relocation, high cross-location
correlation, zero buying cost --- the pre-buying hedge channel should
attenuate or collapse.
Confirmation of these patterns will establish that the welfare gain reported
in Table~\ref{tab:cev_baseline} reflects the model's structural mechanism
rather than a calibration artefact.
\textsc{[Fill in pass/fail outcomes from output/diagnostics/p6\_option1\_decomposition.md.]}

%              Requires amsmath, booktabs, multirow, siunitx, graphicx, caption.

% ─────────────────────────────────────────────────────────────────────────────
\section{Results}
\label{sec:results}
% ─────────────────────────────────────────────────────────────────────────────

We report the baseline welfare results in Section~\ref{sec:results:baseline},
the channel decomposition in Section~\ref{sec:results:decomp}, the
pre-registered falsification tests in Section~\ref{sec:results:falsification},
and sensitivity analysis in Section~\ref{sec:results:sensitivity}.

All welfare gains are expressed as consumption equivalent variations (CEV) at the
initial state
$\mathbf{s}_0 = (t{=}1,\,w_{\text{mid}},\,z_{\text{mid}},\,\ell{=}A,\,x_{A,0}{=}0,\,x_{B,0}{=}0)$:
\begin{equation}
    \text{CEV}(R_a \text{ vs } R_b) \;=\;
    \left(\frac{V_1^{R_a}(\mathbf{s}_0)}{V_1^{R_b}(\mathbf{s}_0)}\right)^{1/(1-\gamma)}
    - 1,
    \label{eq:cev}
\end{equation}
where $V_1^{R}(\mathbf{s}_0)$ is the time-1 value function under regime $R$
evaluated at $\mathbf{s}_0$.%
\footnote{Both value functions are negative under CRRA with finite wealth,
so the ratio is positive and the formula is well-defined for $\gamma \neq 1$.
The CEV measures the fraction of lifetime consumption in regime $R_b$ that the
household would give up to gain permanent access to regime $R_a$.}
Appendix~\ref{app:numerical} reports convergence diagnostics and full-grid
robustness.

% ─────────────────────────────────────────────────────────────────────────────
\subsection{Baseline Welfare Results}
\label{sec:results:baseline}
% ─────────────────────────────────────────────────────────────────────────────

The headline welfare result is reported in Table~\ref{tab:cev_baseline}.
The first row gives $\text{CEV}(\text{E2\_2L vs E1\_2L})$, the lifetime
welfare value of tokenization relative to traditional binary homeownership.
Rows 2--4 decompose this gain into three channels (detailed in
Section~\ref{sec:results:decomp}).
Row 5 reports $\text{CEV}(\text{E0 vs E1\_2L})$ as a sanity-check benchmark:
the welfare cost of pure renting relative to traditional ownership.

\begin{table}[t]
\centering
\caption{Welfare Value of Tokenization: Baseline and Robustness Panel
         (CEV, percentage points of lifetime consumption)}
\label{tab:cev_baseline}
\small
\begin{tabular}{lcccccc}
\toprule
 & Baseline & $\rho_{AB}=0$ & $\rho_{AB}=0.50$ & $\rho_{AB}=0.95$
 & $p_{\text{reloc}}=0$ & $p_{\text{reloc}}=0.12$ \\
\midrule
$\text{CEV}(\text{E2\_2L vs E1\_2L})$
    & \textsc{[P]} & \textsc{[P]} & \textsc{[P]} & \textsc{[P]}
    & \textsc{[P]} & \textsc{[P]} \\[2pt]
\quad Avoided-tx channel
    & \textsc{[P]} & \textsc{[P]} & \textsc{[P]} & \textsc{[P]}
    & \textsc{[P]} & \textsc{[P]} \\
\quad Maintained-hedge channel
    & \textsc{[P]} & \textsc{[P]} & \textsc{[P]} & \textsc{[P]}
    & \textsc{[P]} & \textsc{[P]} \\
\quad Cross-term $\xi$
    & \textsc{[P]} & \textsc{[P]} & \textsc{[P]} & \textsc{[P]}
    & \textsc{[P]} & \textsc{[P]} \\[4pt]
$\text{CEV}(\text{E0 vs E1\_2L})$ (renter benchmark)
    & \textsc{[P]} & --- & --- & --- & --- & --- \\
\addlinespace
Mean $\bar{x}_B$ at $\ell{=}A$
    & \textsc{[P]} & \textsc{[P]} & \textsc{[P]} & \textsc{[P]}
    & \textsc{[P]} & \textsc{[P]} \\
\bottomrule
\end{tabular}
\vspace{4pt}
\begin{minipage}{\linewidth}
\small\textit{Notes:}
Baseline: $\gamma=5$, $\beta=0.96$, $\rho_{AB}=0.50$,
$p_{\text{reloc,work}}=0.06$, $\tau_{\text{sell}}=0.06$,
$\tau_{\text{buy}}=0.025$, $\tau_{\text{token}}=0.01$.
Grid: $N_W=15$, $N_Z=5$, $N_{x_{\text{prev}}}=3$ (coarse; full-grid
robustness at $N_W=40$, $N_Z=9$, $N_{x_{\text{prev}}}=5$ in
Appendix~\ref{app:numerical}).
\textsc{[P]} = placeholder pending server1 baseline runs.
\end{minipage}
\end{table}

The primary diagnostic for mechanism activation is $\bar{x}_B$ at $\ell{=}A$:
the mean portfolio share of location-$B$ tokens held by households residing
at location $A$.
Under E1\_2L, $\bar{x}_B = 0$ by admissibility (binary tenure forbids
non-occupancy-location ownership).
Under E2\_2L, a positive $\bar{x}_B > 0$ at $\ell{=}A$ indicates that
households voluntarily pre-accumulate location-$B$ tokens before relocation ---
the mechanism that is uniquely enabled by token portability and the 6D state
extension (Option~1).
This is the key discriminator from the v3 Option~3 approximation, under which
$\bar{x}_B = 0$ held at all mobility rates.

Figure~\ref{fig:lifecycle_profiles} illustrates the mechanism at the
lifecycle level: the age profiles of mean $x_A$ and $x_B$ holdings
under E1\_2L and E2\_2L.
The key prediction is that $\bar{x}_B$ rises during working years (ages 25--65),
when $p_{\text{reloc,work}} = 0.06$, and declines post-retirement when the
relocation probability falls to $p_{\text{reloc,ret}} = 0.02$.
This pattern --- confirmed or denied by the server1 policy-array output ---
constitutes the strongest visual evidence for the pre-buying hedge mechanism.

\begin{figure}[t]
\centering
%% PLACEHOLDER — replace with lifecycle profile plots once server1 policy arrays land.
%% See paper/exhibit_memos/fig1_lifecycle_profiles.md for full production spec.
\fbox{%
  \begin{minipage}{0.88\linewidth}
  \vspace{2.5cm}
  \centering
  {\sffamily\small [PLACEHOLDER: Figure~1 — Lifecycle token-holding profiles.
  Panel A (E1\_2L): mean $x_A$ (solid) and $x_B = 0$ (dashed) over ages 25--80.
  Panel B (E2\_2L): mean $x_A$ (solid) and mean $x_B > 0$ (dashed) over ages 25--80.
  Source: per-age mean of \texttt{xA\_policy} and \texttt{xB\_policy} arrays at $\ell{=}A$,
  averaged over feasible $(w, z, x_{\text{prev}})$ states.
  Script: \texttt{scripts/plot\_lifecycle\_profiles.py}.]}
  \vspace{2.5cm}
  \end{minipage}%
}
\caption{Lifecycle mean token-holding profiles under traditional ownership (E1\_2L,
         left panel) and tokenized ownership (E2\_2L, right panel) at location $A$.
         E2\_2L households pre-accumulate location-$B$ tokens during working years
         (positive $\bar{x}_B$), saving the buying cost $\tau_{\text{buy}}$ on
         anticipated future relocation.
         E1\_2L households hold $\bar{x}_B = 0$ at all ages by admissibility.
         \textsc{[P]} = placeholder pending server1 run.}
\label{fig:lifecycle_profiles}
\end{figure}

% ─────────────────────────────────────────────────────────────────────────────
\subsection{Channel Decomposition}
\label{sec:results:decomp}
% ─────────────────────────────────────────────────────────────────────────────

We decompose the total welfare gain using three regime runs.
Let E1\_2L$_{\text{NT}}$ denote the no-transaction-cost counterfactual
($\tau_{\text{sell}} = \tau_{\text{buy}} = 0$, binary tenure otherwise).
The decomposition is:
\begin{align}
    \underbrace{\text{CEV}(\text{E2\_2L vs E1\_2L})}_{\text{total}}
    &\;=\;
    \underbrace{\text{CEV}(\text{E1\_2L}_{\text{NT}} \text{ vs E1\_2L})}_{\text{avoided-tx}}
    \;+\;
    \underbrace{\text{CEV}(\text{E2\_2L vs E1\_2L}_{\text{NT}})}_{\text{maintained-hedge}}
    \;+\; \xi,
    \label{eq:decomp}
\end{align}
where $\xi$ is the cross-term.
This decomposition is pre-registered in \texttt{docs/welfare\_decomp\_v4.md}
and is exact up to the cross-term, which was empirically near zero in v3
($\xi = 0.021$\%) and is expected to remain small in v4.

\begin{table}[t]
\centering
\caption{Channel Decomposition of the Tokenization Welfare Gain
         (CEV, percentage points of lifetime consumption)}
\label{tab:decomp}
\small
\begin{tabular}{lccl}
\toprule
Channel & CEV (\%) & Share of total & Description \\
\midrule
Total: $\text{CEV}(\text{E2\_2L vs E1\_2L})$
    & \textsc{[P]} & 100\% & Headline tokenization gain \\
\addlinespace
\quad (i)~Avoided-tx
    & \textsc{[P]} & \textsc{[P]}\%
    & $\tau_{\text{sell}} + \tau_{\text{buy}}$ burden of E1\_2L moves \\
\qquad of which $\tau_{\text{sell}}$
    & \textsc{[P]} & \textsc{[P]}\%
    & NAR selling commission (6\%) \\
\qquad of which $\tau_{\text{buy}}$
    & \textsc{[P]} & \textsc{[P]}\%
    & Buying closing costs (2.5\%) \\[4pt]
\quad (ii)~Maintained-hedge
    & \textsc{[P]} & \textsc{[P]}\%
    & Continuous-x + pre-buying hedge vs E1\_2L$_\text{NT}$ \\[4pt]
\quad (iii)~Cross-term $\xi$
    & \textsc{[P]} & \textsc{[P]}\%
    & Interaction; expected $\approx 0$ \\
\midrule
v3 Option~3 (no state ext.)
    & $+4.255$\% & (reference) & Previous fire's upper-bound approx \\
\quad of which avoided-tx
    & $+0.816$\% & 19.2\% & From v3 full-grid run \\
\quad of which continuous-x
    & $+3.411$\% & 80.2\% & Liu (2021) territory \\
\quad of which pre-buying hedge
    & $\approx 0$\% & $\approx 0$\% & Absent without state ext. (Option~3 limit) \\
\bottomrule
\end{tabular}
\vspace{4pt}
\begin{minipage}{\linewidth}
\small\textit{Notes:}
Upper panel: v4 Option~1 results (\textsc{[P]} = placeholder).
Lower panel: v3 Option~3 reference numbers (confirmed from server1 run,
commit 186da13).
The maintained-hedge channel in v4 encompasses both the continuous-x
subchannel (present in v3) and the new pre-buying hedge subchannel
enabled by the 6D state; a fourth regime decomposition
($\text{E2\_2L}_{\text{NT}}$) can further separate the two if H1 is confirmed.
\end{minipage}
\end{table}

The maintained-hedge channel in v4 captures two sub-mechanisms:
\begin{enumerate}
    \item \emph{Continuous-$x$ sub-channel} (present in v3): continuous
    fractional ownership of the occupied unit avoids the binary kink at
    $x_\ell = 1$, delivering a Liu~(2021)-style indivisibility relaxation.
    \item \emph{Pre-buying hedge sub-channel} (new in v4): by tracking
    $(x_{A,t-1}, x_{B,t-1})$ in the state, households face a per-period
    buying cost $\tau_{\text{buy}} \cdot \max(\Delta x, 0)$ and can
    pre-accumulate location-$B$ tokens at location $A$, arriving at $B$
    with $x_{B,t-1} > 0$ and avoiding the lump-sum purchase cost.
    The expected per-period premium per unit held is
    $p_{\text{reloc}} \times \tau_{\text{buy}} = 0.06 \times 0.025 = 0.15$\%
    (see Appendix~\ref{app:proof}).
\end{enumerate}

% ─────────────────────────────────────────────────────────────────────────────
\subsection{Falsification Tests}
\label{sec:results:falsification}
% ─────────────────────────────────────────────────────────────────────────────

We conduct three pre-registered falsification tests designed to identify
the pre-buying hedge channel.
These tests were pre-registered in \texttt{docs/welfare\_decomp\_v4.md}
before server1 runs were executed.
Crucially, the analogous tests \emph{failed} under v3 Option~3 (which lacked
the state extension): $\bar{x}_B$ remained zero at all mobility rates,
and the CEV was insensitive to $p_{\text{reloc}}$ and $\rho_{AB}$.

\begin{table}[t]
\centering
\caption{Pre-Registered Falsification Tests
         (Deviations from Baseline)}
\label{tab:falsification}
\small
\begin{tabular}{llcccc}
\toprule
Test & Counterfactual & $\text{CEV}$ & $\bar{x}_B$ at $\ell{=}A$
     & Pass criterion & Result \\
\midrule
(r) Relocation disabled
    & $p_{\text{reloc}} = 0$
    & \textsc{[P]}
    & \textsc{[P]}
    & $\bar{x}_B \to 0$;\; $\Delta\text{CEV} < -0.5$pp
    & \textsc{[P]} \\[4pt]
(m) Perfect correlation
    & $\rho_{AB} = 0.95$
    & \textsc{[P]}
    & \textsc{[P]}
    & $\bar{x}_B$ decreases;\; $\Delta\text{CEV} < 0$
    & \textsc{[P]} \\[4pt]
(q) No buying cost
    & $\tau_{\text{buy}} = 0$
    & \textsc{[P]}
    & \textsc{[P]}
    & Hedge channel $\to 0$;\; $\bar{x}_B \to 0$
    & \textsc{[P]} \\
\bottomrule
\end{tabular}
\vspace{4pt}
\begin{minipage}{\linewidth}
\small\textit{Notes:}
All tests compare to the baseline E2\_2L v4 run.
Pass criterion stated prior to server1 runs.
``Hedge channel'' in test (q) is
$\text{CEV}(\text{E2\_2L}_{\tau=0} \text{ vs E1\_2L}_{\text{NT}}) - \text{CEV}(\text{E2\_2L}_{\text{baseline}} \text{ vs E1\_2L}_{\text{NT}})$.
\textsc{[P]} = placeholder.
\end{minipage}
\end{table}

Test (r) shuts off the relocation shock entirely.
Without future relocation, there is no value to pre-accumulating
location-$B$ tokens while at $A$; the pre-buying hedge channel
should collapse.
If test (r) passes --- i.e., $\bar{x}_B \to 0$ when $p_{\text{reloc}} = 0$
--- it provides strong evidence that the positive $\bar{x}_B$ in the
baseline is driven by the pre-buying motive and not by a model artefact.

Test (m) raises the cross-location correlation to $\rho_{AB} = 0.95$,
nearly eliminating diversification benefit.
When location-$A$ and location-$B$ price processes are nearly identical,
holding $x_B$ at $\ell = A$ provides little hedge against diverging location returns.
The CEV and $\bar{x}_B$ should both decrease relative to the baseline.

Test (q) eliminates the buying cost ($\tau_{\text{buy}} = 0$), which
removes the financial incentive to pre-accumulate $x_B$ (there is no lump-sum
arrival cost to avoid).
The pre-buying hedge channel should collapse to zero.

% ─────────────────────────────────────────────────────────────────────────────
\subsection{Sensitivity Analysis}
\label{sec:results:sensitivity}
% ─────────────────────────────────────────────────────────────────────────────

Figure~\ref{fig:sensitivity_heatmap} depicts the total welfare gain
$\text{CEV}(\text{E2\_2L vs E1\_2L})$ as a function of the two primary
mechanism parameters: the cross-location correlation $\rho_{AB}$ and
the relocation probability $p_{\text{reloc}}$.
The full sweep design follows the pre-registered plan in
\texttt{docs/sensitivity\_grid\_v4.md}; only a subset is shown here.

\begin{figure}[t]
\centering
%% PLACEHOLDER — replace with actual heatmap once sweep results available
\fbox{%
  \begin{minipage}{0.72\linewidth}
  \vspace{2.5cm}
  \centering
  {\sffamily\small [PLACEHOLDER: Figure~2 — CEV sensitivity heatmap.
  X-axis: $\rho_{AB} \in \{0, 0.25, 0.50, 0.75, 0.95\}$.
  Y-axis: $p_{\text{reloc}} \in \{0, 0.02, 0.06, 0.12\}$.
  Heat colour: total CEV(\%), from dark (low) to bright (high).
  Source: sweep outputs from \texttt{scripts/sweep\_rhoAB.sh} and
  \texttt{scripts/sweep\_prelocate.sh} once server1 runs complete.]}
  \vspace{2.5cm}
  \end{minipage}%
}
\caption{Welfare sensitivity to cross-location correlation and relocation rate.
         Each cell reports $\text{CEV}(\text{E2\_2L vs E1\_2L})$ at
         $(\rho_{AB},\, p_{\text{reloc}})$.
         The mechanism prediction is a positive gradient in $p_{\text{reloc}}$
         (higher mobility $\Rightarrow$ higher hedge value) and a negative
         gradient in $\rho_{AB}$ (more correlated locations $\Rightarrow$
         smaller diversification benefit).}
\label{fig:sensitivity_heatmap}
\end{figure}

Table~\ref{tab:sensitivity} reports selected cross-sections of the
sensitivity grid for other key parameters.

\begin{table}[t]
\centering
\caption{Sensitivity of Headline CEV to Key Parameters
         (Total CEV, \%; baseline value on diagonal)}
\label{tab:sensitivity}
\small
\begin{tabular}{lcccccc}
\toprule
Parameter & & \multicolumn{5}{c}{Values} \\
\midrule
Risk aversion $\gamma$
    & & 3 & \textbf{5} & 8 & & \\
\quad $\text{CEV}(\text{E2\_2L vs E1\_2L})$
    & & \textsc{[P]} & \textsc{[P]} & \textsc{[P]} & & \\
\addlinespace
Buying cost $\tau_{\text{buy}}$
    & & 0\% & 1\% & \textbf{2.5\%} & 4\% & 6\% \\
\quad $\text{CEV}(\text{E2\_2L vs E1\_2L})$
    & & \textsc{[P]} & \textsc{[P]} & \textsc{[P]} & \textsc{[P]} & \textsc{[P]} \\
\addlinespace
Relocation rate $p_{\text{reloc,work}}$
    & & 0 & 0.02 & \textbf{0.06} & 0.12 & \\
\quad $\text{CEV}(\text{E2\_2L vs E1\_2L})$
    & & \textsc{[P]} & \textsc{[P]} & \textsc{[P]} & \textsc{[P]} & \\
\quad $\bar{x}_B$ at $\ell{=}A$
    & & \textsc{[P]} & \textsc{[P]} & \textsc{[P]} & \textsc{[P]} & \\
\addlinespace
Cross-location corr.\ $\rho_{AB}$
    & & 0 & 0.25 & \textbf{0.50} & 0.75 & 0.95 \\
\quad $\text{CEV}(\text{E2\_2L vs E1\_2L})$
    & & \textsc{[P]} & \textsc{[P]} & \textsc{[P]} & \textsc{[P]} & \textsc{[P]} \\
\addlinespace
\multicolumn{7}{l}{\textit{Asymmetric robustness (solver v4 extension, fire 22)}} \\
Location-$B$ return premium $\Delta\mu_B$
    & & $-1\%$ & \textbf{0\%} & $+1\%$ & & \\
\quad $\text{CEV}(\text{E2\_2L vs E1\_2L})$
    & & \textsc{[P]} & \textsc{[P]} & \textsc{[P]} & & \\
\quad $\bar{x}_B$ at $\ell{=}A$
    & & \textsc{[P]} & \textsc{[P]} & \textsc{[P]} & & \\[2pt]
Directional mobility ($p_{A\to B}$, $p_{B\to A}$)
    & & (0.06, 0.06) & (0.10, 0.03) & (0.10, 0.10) & & \\
\quad $\text{CEV}(\text{E2\_2L vs E1\_2L})$
    & & \textsc{[P]} & \textsc{[P]} & \textsc{[P]} & & \\
\addlinespace
\multicolumn{7}{l}{\textit{Mortgage activation (LTV constraint)}} \\
LTV ceiling $\theta_{\max}$
    & & \textbf{0} & 0.50 & 0.80 & & \\
\quad $\text{CEV}(\text{E2\_2L vs E1\_2L})$
    & & \textsc{[P]} & \textsc{[P]} & \textsc{[P]} & & \\
\bottomrule
\end{tabular}
\vspace{4pt}
\begin{minipage}{\linewidth}
\small\textit{Notes:}
Bold entries mark the baseline value.
One parameter varies at a time; all others held at baseline.
Sweep scripts: \texttt{scripts/sweep\_rhoAB.sh},
\texttt{scripts/sweep\_prelocate.sh}, \texttt{scripts/sweep\_txcost.sh},
\texttt{scripts/sweep\_asymmetric.sh} (location-$B$ return premium and directional mobility),
\texttt{scripts/sweep\_mortgage.sh} (LTV).
$\Delta\mu_B = \mu_{h,B} - \mu_{h,A}$: the excess Jensen-corrected log-mean return of
location $B$ relative to $A$.
Positive $\Delta\mu_B$ strengthens the pre-buying motive for $x_B$ by adding a return
premium to the tau\_buy-saving incentive.
\textsc{[P]} = placeholder pending server1 sweep runs.
\end{minipage}
\end{table}

Two monotonicity predictions follow from the model mechanism.
First, for $\tau_{\text{buy}}$: as the buying cost rises, the pre-accumulation
motive strengthens and $\bar{x}_B$ should increase, lifting the hedge sub-channel.
Second, for $p_{\text{reloc}}$: higher mobility raises the expected value of
pre-accumulated $x_B$ holdings; at $p_{\text{reloc}} = 0$ the hedge channel
collapses (falsification test r).
If these monotonicity patterns hold in the sweep results, they constitute
additional evidence for the pre-buying hedge mechanism.

% ─────────────────────────────────────────────────────────────────────────────
\subsection{Summary of Mechanism Evidence}
\label{sec:results:summary}
% ─────────────────────────────────────────────────────────────────────────────

The results establish \textsc{[P]} as the headline lifetime welfare gain from
tokenization under the baseline calibration.
The gain decomposes additively into an avoided-transaction-cost channel
(\textsc{[P]}\%) and a maintained-hedge channel (\textsc{[P]}\%), with
a near-zero cross-term.
The maintained-hedge channel is itself divisible into the continuous-$x$
sub-channel present in prior work (Liu, 2021) and the novel pre-buying hedge
sub-channel unlocked by the 6D state extension.

The three falsification tests provide mechanism identification.
Under all three counterfactuals --- zero relocation, high cross-location
correlation, zero buying cost --- the pre-buying hedge channel should
attenuate or collapse.
Confirmation of these patterns will establish that the welfare gain reported
in Table~\ref{tab:cev_baseline} reflects the model's structural mechanism
rather than a calibration artefact.
\textsc{[Fill in pass/fail outcomes from output/diagnostics/p6\_option1\_decomposition.md.]}

%              Requires amsmath, booktabs, multirow, siunitx, graphicx, caption.

% ─────────────────────────────────────────────────────────────────────────────
\section{Results}
\label{sec:results}
% ─────────────────────────────────────────────────────────────────────────────

We report the baseline welfare results in Section~\ref{sec:results:baseline},
the channel decomposition in Section~\ref{sec:results:decomp}, the
pre-registered falsification tests in Section~\ref{sec:results:falsification},
and sensitivity analysis in Section~\ref{sec:results:sensitivity}.

All welfare gains are expressed as consumption equivalent variations (CEV) at the
initial state
$\mathbf{s}_0 = (t{=}1,\,w_{\text{mid}},\,z_{\text{mid}},\,\ell{=}A,\,x_{A,0}{=}0,\,x_{B,0}{=}0)$:
\begin{equation}
    \text{CEV}(R_a \text{ vs } R_b) \;=\;
    \left(\frac{V_1^{R_a}(\mathbf{s}_0)}{V_1^{R_b}(\mathbf{s}_0)}\right)^{1/(1-\gamma)}
    - 1,
    \label{eq:cev}
\end{equation}
where $V_1^{R}(\mathbf{s}_0)$ is the time-1 value function under regime $R$
evaluated at $\mathbf{s}_0$.%
\footnote{Both value functions are negative under CRRA with finite wealth,
so the ratio is positive and the formula is well-defined for $\gamma \neq 1$.
The CEV measures the fraction of lifetime consumption in regime $R_b$ that the
household would give up to gain permanent access to regime $R_a$.}
Appendix~\ref{app:numerical} reports convergence diagnostics and full-grid
robustness.

% ─────────────────────────────────────────────────────────────────────────────
\subsection{Baseline Welfare Results}
\label{sec:results:baseline}
% ─────────────────────────────────────────────────────────────────────────────

The headline welfare result is reported in Table~\ref{tab:cev_baseline}.
The first row gives $\text{CEV}(\text{E2\_2L vs E1\_2L})$, the lifetime
welfare value of tokenization relative to traditional binary homeownership.
Rows 2--4 decompose this gain into three channels (detailed in
Section~\ref{sec:results:decomp}).
Row 5 reports $\text{CEV}(\text{E0 vs E1\_2L})$ as a sanity-check benchmark:
the welfare cost of pure renting relative to traditional ownership.

\begin{table}[t]
\centering
\caption{Welfare Value of Tokenization: Baseline and Robustness Panel
         (CEV, percentage points of lifetime consumption)}
\label{tab:cev_baseline}
\small
\begin{tabular}{lcccccc}
\toprule
 & Baseline & $\rho_{AB}=0$ & $\rho_{AB}=0.50$ & $\rho_{AB}=0.95$
 & $p_{\text{reloc}}=0$ & $p_{\text{reloc}}=0.12$ \\
\midrule
$\text{CEV}(\text{E2\_2L vs E1\_2L})$
    & \textsc{[P]} & \textsc{[P]} & \textsc{[P]} & \textsc{[P]}
    & \textsc{[P]} & \textsc{[P]} \\[2pt]
\quad Avoided-tx channel
    & \textsc{[P]} & \textsc{[P]} & \textsc{[P]} & \textsc{[P]}
    & \textsc{[P]} & \textsc{[P]} \\
\quad Maintained-hedge channel
    & \textsc{[P]} & \textsc{[P]} & \textsc{[P]} & \textsc{[P]}
    & \textsc{[P]} & \textsc{[P]} \\
\quad Cross-term $\xi$
    & \textsc{[P]} & \textsc{[P]} & \textsc{[P]} & \textsc{[P]}
    & \textsc{[P]} & \textsc{[P]} \\[4pt]
$\text{CEV}(\text{E0 vs E1\_2L})$ (renter benchmark)
    & \textsc{[P]} & --- & --- & --- & --- & --- \\
\addlinespace
Mean $\bar{x}_B$ at $\ell{=}A$
    & \textsc{[P]} & \textsc{[P]} & \textsc{[P]} & \textsc{[P]}
    & \textsc{[P]} & \textsc{[P]} \\
\bottomrule
\end{tabular}
\vspace{4pt}
\begin{minipage}{\linewidth}
\small\textit{Notes:}
Baseline: $\gamma=5$, $\beta=0.96$, $\rho_{AB}=0.50$,
$p_{\text{reloc,work}}=0.06$, $\tau_{\text{sell}}=0.06$,
$\tau_{\text{buy}}=0.025$, $\tau_{\text{token}}=0.01$.
Grid: $N_W=15$, $N_Z=5$, $N_{x_{\text{prev}}}=3$ (coarse; full-grid
robustness at $N_W=40$, $N_Z=9$, $N_{x_{\text{prev}}}=5$ in
Appendix~\ref{app:numerical}).
\textsc{[P]} = placeholder pending server1 baseline runs.
\end{minipage}
\end{table}

The primary diagnostic for mechanism activation is $\bar{x}_B$ at $\ell{=}A$:
the mean portfolio share of location-$B$ tokens held by households residing
at location $A$.
Under E1\_2L, $\bar{x}_B = 0$ by admissibility (binary tenure forbids
non-occupancy-location ownership).
Under E2\_2L, a positive $\bar{x}_B > 0$ at $\ell{=}A$ indicates that
households voluntarily pre-accumulate location-$B$ tokens before relocation ---
the mechanism that is uniquely enabled by token portability and the 6D state
extension (Option~1).
This is the key discriminator from the v3 Option~3 approximation, under which
$\bar{x}_B = 0$ held at all mobility rates.

Figure~\ref{fig:lifecycle_profiles} illustrates the mechanism at the
lifecycle level: the age profiles of mean $x_A$ and $x_B$ holdings
under E1\_2L and E2\_2L.
The key prediction is that $\bar{x}_B$ rises during working years (ages 25--65),
when $p_{\text{reloc,work}} = 0.06$, and declines post-retirement when the
relocation probability falls to $p_{\text{reloc,ret}} = 0.02$.
This pattern --- confirmed or denied by the server1 policy-array output ---
constitutes the strongest visual evidence for the pre-buying hedge mechanism.

\begin{figure}[t]
\centering
%% PLACEHOLDER — replace with lifecycle profile plots once server1 policy arrays land.
%% See paper/exhibit_memos/fig1_lifecycle_profiles.md for full production spec.
\fbox{%
  \begin{minipage}{0.88\linewidth}
  \vspace{2.5cm}
  \centering
  {\sffamily\small [PLACEHOLDER: Figure~1 — Lifecycle token-holding profiles.
  Panel A (E1\_2L): mean $x_A$ (solid) and $x_B = 0$ (dashed) over ages 25--80.
  Panel B (E2\_2L): mean $x_A$ (solid) and mean $x_B > 0$ (dashed) over ages 25--80.
  Source: per-age mean of \texttt{xA\_policy} and \texttt{xB\_policy} arrays at $\ell{=}A$,
  averaged over feasible $(w, z, x_{\text{prev}})$ states.
  Script: \texttt{scripts/plot\_lifecycle\_profiles.py}.]}
  \vspace{2.5cm}
  \end{minipage}%
}
\caption{Lifecycle mean token-holding profiles under traditional ownership (E1\_2L,
         left panel) and tokenized ownership (E2\_2L, right panel) at location $A$.
         E2\_2L households pre-accumulate location-$B$ tokens during working years
         (positive $\bar{x}_B$), saving the buying cost $\tau_{\text{buy}}$ on
         anticipated future relocation.
         E1\_2L households hold $\bar{x}_B = 0$ at all ages by admissibility.
         \textsc{[P]} = placeholder pending server1 run.}
\label{fig:lifecycle_profiles}
\end{figure}

% ─────────────────────────────────────────────────────────────────────────────
\subsection{Channel Decomposition}
\label{sec:results:decomp}
% ─────────────────────────────────────────────────────────────────────────────

We decompose the total welfare gain using three regime runs.
Let E1\_2L$_{\text{NT}}$ denote the no-transaction-cost counterfactual
($\tau_{\text{sell}} = \tau_{\text{buy}} = 0$, binary tenure otherwise).
The decomposition is:
\begin{align}
    \underbrace{\text{CEV}(\text{E2\_2L vs E1\_2L})}_{\text{total}}
    &\;=\;
    \underbrace{\text{CEV}(\text{E1\_2L}_{\text{NT}} \text{ vs E1\_2L})}_{\text{avoided-tx}}
    \;+\;
    \underbrace{\text{CEV}(\text{E2\_2L vs E1\_2L}_{\text{NT}})}_{\text{maintained-hedge}}
    \;+\; \xi,
    \label{eq:decomp}
\end{align}
where $\xi$ is the cross-term.
This decomposition is pre-registered in \texttt{docs/welfare\_decomp\_v4.md}
and is exact up to the cross-term, which was empirically near zero in v3
($\xi = 0.021$\%) and is expected to remain small in v4.

\begin{table}[t]
\centering
\caption{Channel Decomposition of the Tokenization Welfare Gain
         (CEV, percentage points of lifetime consumption)}
\label{tab:decomp}
\small
\begin{tabular}{lccl}
\toprule
Channel & CEV (\%) & Share of total & Description \\
\midrule
Total: $\text{CEV}(\text{E2\_2L vs E1\_2L})$
    & \textsc{[P]} & 100\% & Headline tokenization gain \\
\addlinespace
\quad (i)~Avoided-tx
    & \textsc{[P]} & \textsc{[P]}\%
    & $\tau_{\text{sell}} + \tau_{\text{buy}}$ burden of E1\_2L moves \\
\qquad of which $\tau_{\text{sell}}$
    & \textsc{[P]} & \textsc{[P]}\%
    & NAR selling commission (6\%) \\
\qquad of which $\tau_{\text{buy}}$
    & \textsc{[P]} & \textsc{[P]}\%
    & Buying closing costs (2.5\%) \\[4pt]
\quad (ii)~Maintained-hedge
    & \textsc{[P]} & \textsc{[P]}\%
    & Continuous-x + pre-buying hedge vs E1\_2L$_\text{NT}$ \\[4pt]
\quad (iii)~Cross-term $\xi$
    & \textsc{[P]} & \textsc{[P]}\%
    & Interaction; expected $\approx 0$ \\
\midrule
v3 Option~3 (no state ext.)
    & $+4.255$\% & (reference) & Previous fire's upper-bound approx \\
\quad of which avoided-tx
    & $+0.816$\% & 19.2\% & From v3 full-grid run \\
\quad of which continuous-x
    & $+3.411$\% & 80.2\% & Liu (2021) territory \\
\quad of which pre-buying hedge
    & $\approx 0$\% & $\approx 0$\% & Absent without state ext. (Option~3 limit) \\
\bottomrule
\end{tabular}
\vspace{4pt}
\begin{minipage}{\linewidth}
\small\textit{Notes:}
Upper panel: v4 Option~1 results (\textsc{[P]} = placeholder).
Lower panel: v3 Option~3 reference numbers (confirmed from server1 run,
commit 186da13).
The maintained-hedge channel in v4 encompasses both the continuous-x
subchannel (present in v3) and the new pre-buying hedge subchannel
enabled by the 6D state; a fourth regime decomposition
($\text{E2\_2L}_{\text{NT}}$) can further separate the two if H1 is confirmed.
\end{minipage}
\end{table}

The maintained-hedge channel in v4 captures two sub-mechanisms:
\begin{enumerate}
    \item \emph{Continuous-$x$ sub-channel} (present in v3): continuous
    fractional ownership of the occupied unit avoids the binary kink at
    $x_\ell = 1$, delivering a Liu~(2021)-style indivisibility relaxation.
    \item \emph{Pre-buying hedge sub-channel} (new in v4): by tracking
    $(x_{A,t-1}, x_{B,t-1})$ in the state, households face a per-period
    buying cost $\tau_{\text{buy}} \cdot \max(\Delta x, 0)$ and can
    pre-accumulate location-$B$ tokens at location $A$, arriving at $B$
    with $x_{B,t-1} > 0$ and avoiding the lump-sum purchase cost.
    The expected per-period premium per unit held is
    $p_{\text{reloc}} \times \tau_{\text{buy}} = 0.06 \times 0.025 = 0.15$\%
    (see Appendix~\ref{app:proof}).
\end{enumerate}

% ─────────────────────────────────────────────────────────────────────────────
\subsection{Falsification Tests}
\label{sec:results:falsification}
% ─────────────────────────────────────────────────────────────────────────────

We conduct three pre-registered falsification tests designed to identify
the pre-buying hedge channel.
These tests were pre-registered in \texttt{docs/welfare\_decomp\_v4.md}
before server1 runs were executed.
Crucially, the analogous tests \emph{failed} under v3 Option~3 (which lacked
the state extension): $\bar{x}_B$ remained zero at all mobility rates,
and the CEV was insensitive to $p_{\text{reloc}}$ and $\rho_{AB}$.

\begin{table}[t]
\centering
\caption{Pre-Registered Falsification Tests
         (Deviations from Baseline)}
\label{tab:falsification}
\small
\begin{tabular}{llcccc}
\toprule
Test & Counterfactual & $\text{CEV}$ & $\bar{x}_B$ at $\ell{=}A$
     & Pass criterion & Result \\
\midrule
(r) Relocation disabled
    & $p_{\text{reloc}} = 0$
    & \textsc{[P]}
    & \textsc{[P]}
    & $\bar{x}_B \to 0$;\; $\Delta\text{CEV} < -0.5$pp
    & \textsc{[P]} \\[4pt]
(m) Perfect correlation
    & $\rho_{AB} = 0.95$
    & \textsc{[P]}
    & \textsc{[P]}
    & $\bar{x}_B$ decreases;\; $\Delta\text{CEV} < 0$
    & \textsc{[P]} \\[4pt]
(q) No buying cost
    & $\tau_{\text{buy}} = 0$
    & \textsc{[P]}
    & \textsc{[P]}
    & Hedge channel $\to 0$;\; $\bar{x}_B \to 0$
    & \textsc{[P]} \\
\bottomrule
\end{tabular}
\vspace{4pt}
\begin{minipage}{\linewidth}
\small\textit{Notes:}
All tests compare to the baseline E2\_2L v4 run.
Pass criterion stated prior to server1 runs.
``Hedge channel'' in test (q) is
$\text{CEV}(\text{E2\_2L}_{\tau=0} \text{ vs E1\_2L}_{\text{NT}}) - \text{CEV}(\text{E2\_2L}_{\text{baseline}} \text{ vs E1\_2L}_{\text{NT}})$.
\textsc{[P]} = placeholder.
\end{minipage}
\end{table}

Test (r) shuts off the relocation shock entirely.
Without future relocation, there is no value to pre-accumulating
location-$B$ tokens while at $A$; the pre-buying hedge channel
should collapse.
If test (r) passes --- i.e., $\bar{x}_B \to 0$ when $p_{\text{reloc}} = 0$
--- it provides strong evidence that the positive $\bar{x}_B$ in the
baseline is driven by the pre-buying motive and not by a model artefact.

Test (m) raises the cross-location correlation to $\rho_{AB} = 0.95$,
nearly eliminating diversification benefit.
When location-$A$ and location-$B$ price processes are nearly identical,
holding $x_B$ at $\ell = A$ provides little hedge against diverging location returns.
The CEV and $\bar{x}_B$ should both decrease relative to the baseline.

Test (q) eliminates the buying cost ($\tau_{\text{buy}} = 0$), which
removes the financial incentive to pre-accumulate $x_B$ (there is no lump-sum
arrival cost to avoid).
The pre-buying hedge channel should collapse to zero.

% ─────────────────────────────────────────────────────────────────────────────
\subsection{Sensitivity Analysis}
\label{sec:results:sensitivity}
% ─────────────────────────────────────────────────────────────────────────────

Figure~\ref{fig:sensitivity_heatmap} depicts the total welfare gain
$\text{CEV}(\text{E2\_2L vs E1\_2L})$ as a function of the two primary
mechanism parameters: the cross-location correlation $\rho_{AB}$ and
the relocation probability $p_{\text{reloc}}$.
The full sweep design follows the pre-registered plan in
\texttt{docs/sensitivity\_grid\_v4.md}; only a subset is shown here.

\begin{figure}[t]
\centering
%% PLACEHOLDER — replace with actual heatmap once sweep results available
\fbox{%
  \begin{minipage}{0.72\linewidth}
  \vspace{2.5cm}
  \centering
  {\sffamily\small [PLACEHOLDER: Figure~2 — CEV sensitivity heatmap.
  X-axis: $\rho_{AB} \in \{0, 0.25, 0.50, 0.75, 0.95\}$.
  Y-axis: $p_{\text{reloc}} \in \{0, 0.02, 0.06, 0.12\}$.
  Heat colour: total CEV(\%), from dark (low) to bright (high).
  Source: sweep outputs from \texttt{scripts/sweep\_rhoAB.sh} and
  \texttt{scripts/sweep\_prelocate.sh} once server1 runs complete.]}
  \vspace{2.5cm}
  \end{minipage}%
}
\caption{Welfare sensitivity to cross-location correlation and relocation rate.
         Each cell reports $\text{CEV}(\text{E2\_2L vs E1\_2L})$ at
         $(\rho_{AB},\, p_{\text{reloc}})$.
         The mechanism prediction is a positive gradient in $p_{\text{reloc}}$
         (higher mobility $\Rightarrow$ higher hedge value) and a negative
         gradient in $\rho_{AB}$ (more correlated locations $\Rightarrow$
         smaller diversification benefit).}
\label{fig:sensitivity_heatmap}
\end{figure}

Table~\ref{tab:sensitivity} reports selected cross-sections of the
sensitivity grid for other key parameters.

\begin{table}[t]
\centering
\caption{Sensitivity of Headline CEV to Key Parameters
         (Total CEV, \%; baseline value on diagonal)}
\label{tab:sensitivity}
\small
\begin{tabular}{lcccccc}
\toprule
Parameter & & \multicolumn{5}{c}{Values} \\
\midrule
Risk aversion $\gamma$
    & & 3 & \textbf{5} & 8 & & \\
\quad $\text{CEV}(\text{E2\_2L vs E1\_2L})$
    & & \textsc{[P]} & \textsc{[P]} & \textsc{[P]} & & \\
\addlinespace
Buying cost $\tau_{\text{buy}}$
    & & 0\% & 1\% & \textbf{2.5\%} & 4\% & 6\% \\
\quad $\text{CEV}(\text{E2\_2L vs E1\_2L})$
    & & \textsc{[P]} & \textsc{[P]} & \textsc{[P]} & \textsc{[P]} & \textsc{[P]} \\
\addlinespace
Relocation rate $p_{\text{reloc,work}}$
    & & 0 & 0.02 & \textbf{0.06} & 0.12 & \\
\quad $\text{CEV}(\text{E2\_2L vs E1\_2L})$
    & & \textsc{[P]} & \textsc{[P]} & \textsc{[P]} & \textsc{[P]} & \\
\quad $\bar{x}_B$ at $\ell{=}A$
    & & \textsc{[P]} & \textsc{[P]} & \textsc{[P]} & \textsc{[P]} & \\
\addlinespace
Cross-location corr.\ $\rho_{AB}$
    & & 0 & 0.25 & \textbf{0.50} & 0.75 & 0.95 \\
\quad $\text{CEV}(\text{E2\_2L vs E1\_2L})$
    & & \textsc{[P]} & \textsc{[P]} & \textsc{[P]} & \textsc{[P]} & \textsc{[P]} \\
\addlinespace
\multicolumn{7}{l}{\textit{Asymmetric robustness (solver v4 extension, fire 22)}} \\
Location-$B$ return premium $\Delta\mu_B$
    & & $-1\%$ & \textbf{0\%} & $+1\%$ & & \\
\quad $\text{CEV}(\text{E2\_2L vs E1\_2L})$
    & & \textsc{[P]} & \textsc{[P]} & \textsc{[P]} & & \\
\quad $\bar{x}_B$ at $\ell{=}A$
    & & \textsc{[P]} & \textsc{[P]} & \textsc{[P]} & & \\[2pt]
Directional mobility ($p_{A\to B}$, $p_{B\to A}$)
    & & (0.06, 0.06) & (0.10, 0.03) & (0.10, 0.10) & & \\
\quad $\text{CEV}(\text{E2\_2L vs E1\_2L})$
    & & \textsc{[P]} & \textsc{[P]} & \textsc{[P]} & & \\
\addlinespace
\multicolumn{7}{l}{\textit{Mortgage activation (LTV constraint)}} \\
LTV ceiling $\theta_{\max}$
    & & \textbf{0} & 0.50 & 0.80 & & \\
\quad $\text{CEV}(\text{E2\_2L vs E1\_2L})$
    & & \textsc{[P]} & \textsc{[P]} & \textsc{[P]} & & \\
\bottomrule
\end{tabular}
\vspace{4pt}
\begin{minipage}{\linewidth}
\small\textit{Notes:}
Bold entries mark the baseline value.
One parameter varies at a time; all others held at baseline.
Sweep scripts: \texttt{scripts/sweep\_rhoAB.sh},
\texttt{scripts/sweep\_prelocate.sh}, \texttt{scripts/sweep\_txcost.sh},
\texttt{scripts/sweep\_asymmetric.sh} (location-$B$ return premium and directional mobility),
\texttt{scripts/sweep\_mortgage.sh} (LTV).
$\Delta\mu_B = \mu_{h,B} - \mu_{h,A}$: the excess Jensen-corrected log-mean return of
location $B$ relative to $A$.
Positive $\Delta\mu_B$ strengthens the pre-buying motive for $x_B$ by adding a return
premium to the tau\_buy-saving incentive.
\textsc{[P]} = placeholder pending server1 sweep runs.
\end{minipage}
\end{table}

Two monotonicity predictions follow from the model mechanism.
First, for $\tau_{\text{buy}}$: as the buying cost rises, the pre-accumulation
motive strengthens and $\bar{x}_B$ should increase, lifting the hedge sub-channel.
Second, for $p_{\text{reloc}}$: higher mobility raises the expected value of
pre-accumulated $x_B$ holdings; at $p_{\text{reloc}} = 0$ the hedge channel
collapses (falsification test r).
If these monotonicity patterns hold in the sweep results, they constitute
additional evidence for the pre-buying hedge mechanism.

% ─────────────────────────────────────────────────────────────────────────────
\subsection{Summary of Mechanism Evidence}
\label{sec:results:summary}
% ─────────────────────────────────────────────────────────────────────────────

The results establish \textsc{[P]} as the headline lifetime welfare gain from
tokenization under the baseline calibration.
The gain decomposes additively into an avoided-transaction-cost channel
(\textsc{[P]}\%) and a maintained-hedge channel (\textsc{[P]}\%), with
a near-zero cross-term.
The maintained-hedge channel is itself divisible into the continuous-$x$
sub-channel present in prior work (Liu, 2021) and the novel pre-buying hedge
sub-channel unlocked by the 6D state extension.

The three falsification tests provide mechanism identification.
Under all three counterfactuals --- zero relocation, high cross-location
correlation, zero buying cost --- the pre-buying hedge channel should
attenuate or collapse.
Confirmation of these patterns will establish that the welfare gain reported
in Table~\ref{tab:cev_baseline} reflects the model's structural mechanism
rather than a calibration artefact.
\textsc{[Fill in pass/fail outcomes from output/diagnostics/p6\_option1\_decomposition.md.]}

%              Requires amsmath, booktabs, multirow, siunitx, graphicx, caption.

% ─────────────────────────────────────────────────────────────────────────────
\section{Results}
\label{sec:results}
% ─────────────────────────────────────────────────────────────────────────────

We report the baseline welfare results in Section~\ref{sec:results:baseline},
the channel decomposition in Section~\ref{sec:results:decomp}, the
pre-registered falsification tests in Section~\ref{sec:results:falsification},
and sensitivity analysis in Section~\ref{sec:results:sensitivity}.

All welfare gains are expressed as consumption equivalent variations (CEV) at the
initial state
$\mathbf{s}_0 = (t{=}1,\,w_{\text{mid}},\,z_{\text{mid}},\,\ell{=}A,\,x_{A,0}{=}0,\,x_{B,0}{=}0)$:
\begin{equation}
    \text{CEV}(R_a \text{ vs } R_b) \;=\;
    \left(\frac{V_1^{R_a}(\mathbf{s}_0)}{V_1^{R_b}(\mathbf{s}_0)}\right)^{1/(1-\gamma)}
    - 1,
    \label{eq:cev}
\end{equation}
where $V_1^{R}(\mathbf{s}_0)$ is the time-1 value function under regime $R$
evaluated at $\mathbf{s}_0$.%
\footnote{Both value functions are negative under CRRA with finite wealth,
so the ratio is positive and the formula is well-defined for $\gamma \neq 1$.
The CEV measures the fraction of lifetime consumption in regime $R_b$ that the
household would give up to gain permanent access to regime $R_a$.}
Appendix~\ref{app:numerical} reports convergence diagnostics and full-grid
robustness.

% ─────────────────────────────────────────────────────────────────────────────
\subsection{Baseline Welfare Results}
\label{sec:results:baseline}
% ─────────────────────────────────────────────────────────────────────────────

The headline welfare result is reported in Table~\ref{tab:cev_baseline}.
The first row gives $\text{CEV}(\text{E2\_2L vs E1\_2L})$, the lifetime
welfare value of tokenization relative to traditional binary homeownership.
Rows 2--4 decompose this gain into three channels (detailed in
Section~\ref{sec:results:decomp}).
Row 5 reports $\text{CEV}(\text{E0 vs E1\_2L})$ as a sanity-check benchmark:
the welfare cost of pure renting relative to traditional ownership.

\begin{table}[t]
\centering
\caption{Welfare Value of Tokenization: Baseline and Robustness Panel
         (CEV, percentage points of lifetime consumption)}
\label{tab:cev_baseline}
\small
\begin{tabular}{lcccccc}
\toprule
 & Baseline & $\rho_{AB}=0$ & $\rho_{AB}=0.50$ & $\rho_{AB}=0.95$
 & $p_{\text{reloc}}=0$ & $p_{\text{reloc}}=0.12$ \\
\midrule
$\text{CEV}(\text{E2\_2L vs E1\_2L})$
    & \textsc{[P]} & \textsc{[P]} & \textsc{[P]} & \textsc{[P]}
    & \textsc{[P]} & \textsc{[P]} \\[2pt]
\quad Avoided-tx channel
    & \textsc{[P]} & \textsc{[P]} & \textsc{[P]} & \textsc{[P]}
    & \textsc{[P]} & \textsc{[P]} \\
\quad Maintained-hedge channel
    & \textsc{[P]} & \textsc{[P]} & \textsc{[P]} & \textsc{[P]}
    & \textsc{[P]} & \textsc{[P]} \\
\quad Cross-term $\xi$
    & \textsc{[P]} & \textsc{[P]} & \textsc{[P]} & \textsc{[P]}
    & \textsc{[P]} & \textsc{[P]} \\[4pt]
$\text{CEV}(\text{E0 vs E1\_2L})$ (renter benchmark)
    & \textsc{[P]} & --- & --- & --- & --- & --- \\
\addlinespace
Mean $\bar{x}_B$ at $\ell{=}A$
    & \textsc{[P]} & \textsc{[P]} & \textsc{[P]} & \textsc{[P]}
    & \textsc{[P]} & \textsc{[P]} \\
\bottomrule
\end{tabular}
\vspace{4pt}
\begin{minipage}{\linewidth}
\small\textit{Notes:}
Baseline: $\gamma=5$, $\beta=0.96$, $\rho_{AB}=0.50$,
$p_{\text{reloc,work}}=0.06$, $\tau_{\text{sell}}=0.06$,
$\tau_{\text{buy}}=0.025$, $\tau_{\text{token}}=0.01$.
Grid: $N_W=15$, $N_Z=5$, $N_{x_{\text{prev}}}=3$ (coarse; full-grid
robustness at $N_W=40$, $N_Z=9$, $N_{x_{\text{prev}}}=5$ in
Appendix~\ref{app:numerical}).
\textsc{[P]} = placeholder pending server1 baseline runs.
\end{minipage}
\end{table}

The primary diagnostic for mechanism activation is $\bar{x}_B$ at $\ell{=}A$:
the mean portfolio share of location-$B$ tokens held by households residing
at location $A$.
Under E1\_2L, $\bar{x}_B = 0$ by admissibility (binary tenure forbids
non-occupancy-location ownership).
Under E2\_2L, a positive $\bar{x}_B > 0$ at $\ell{=}A$ indicates that
households voluntarily pre-accumulate location-$B$ tokens before relocation ---
the mechanism that is uniquely enabled by token portability and the 6D state
extension (Option~1).
This is the key discriminator from the v3 Option~3 approximation, under which
$\bar{x}_B = 0$ held at all mobility rates.

Figure~\ref{fig:lifecycle_profiles} illustrates the mechanism at the
lifecycle level: the age profiles of mean $x_A$ and $x_B$ holdings
under E1\_2L and E2\_2L.
The key prediction is that $\bar{x}_B$ rises during working years (ages 25--65),
when $p_{\text{reloc,work}} = 0.06$, and declines post-retirement when the
relocation probability falls to $p_{\text{reloc,ret}} = 0.02$.
This pattern --- confirmed or denied by the server1 policy-array output ---
constitutes the strongest visual evidence for the pre-buying hedge mechanism.

\begin{figure}[t]
\centering
%% PLACEHOLDER — replace with lifecycle profile plots once server1 policy arrays land.
%% See paper/exhibit_memos/fig1_lifecycle_profiles.md for full production spec.
\fbox{%
  \begin{minipage}{0.88\linewidth}
  \vspace{2.5cm}
  \centering
  {\sffamily\small [PLACEHOLDER: Figure~1 — Lifecycle token-holding profiles.
  Panel A (E1\_2L): mean $x_A$ (solid) and $x_B = 0$ (dashed) over ages 25--80.
  Panel B (E2\_2L): mean $x_A$ (solid) and mean $x_B > 0$ (dashed) over ages 25--80.
  Source: per-age mean of \texttt{xA\_policy} and \texttt{xB\_policy} arrays at $\ell{=}A$,
  averaged over feasible $(w, z, x_{\text{prev}})$ states.
  Script: \texttt{scripts/plot\_lifecycle\_profiles.py}.]}
  \vspace{2.5cm}
  \end{minipage}%
}
\caption{Lifecycle mean token-holding profiles under traditional ownership (E1\_2L,
         left panel) and tokenized ownership (E2\_2L, right panel) at location $A$.
         E2\_2L households pre-accumulate location-$B$ tokens during working years
         (positive $\bar{x}_B$), saving the buying cost $\tau_{\text{buy}}$ on
         anticipated future relocation.
         E1\_2L households hold $\bar{x}_B = 0$ at all ages by admissibility.
         \textsc{[P]} = placeholder pending server1 run.}
\label{fig:lifecycle_profiles}
\end{figure}

% ─────────────────────────────────────────────────────────────────────────────
\subsection{Channel Decomposition}
\label{sec:results:decomp}
% ─────────────────────────────────────────────────────────────────────────────

We decompose the total welfare gain using three regime runs.
Let E1\_2L$_{\text{NT}}$ denote the no-transaction-cost counterfactual
($\tau_{\text{sell}} = \tau_{\text{buy}} = 0$, binary tenure otherwise).
The decomposition is:
\begin{align}
    \underbrace{\text{CEV}(\text{E2\_2L vs E1\_2L})}_{\text{total}}
    &\;=\;
    \underbrace{\text{CEV}(\text{E1\_2L}_{\text{NT}} \text{ vs E1\_2L})}_{\text{avoided-tx}}
    \;+\;
    \underbrace{\text{CEV}(\text{E2\_2L vs E1\_2L}_{\text{NT}})}_{\text{maintained-hedge}}
    \;+\; \xi,
    \label{eq:decomp}
\end{align}
where $\xi$ is the cross-term.
This decomposition is pre-registered in \texttt{docs/welfare\_decomp\_v4.md}
and is exact up to the cross-term, which was empirically near zero in v3
($\xi = 0.021$\%) and is expected to remain small in v4.

\begin{table}[t]
\centering
\caption{Channel Decomposition of the Tokenization Welfare Gain
         (CEV, percentage points of lifetime consumption)}
\label{tab:decomp}
\small
\begin{tabular}{lccl}
\toprule
Channel & CEV (\%) & Share of total & Description \\
\midrule
Total: $\text{CEV}(\text{E2\_2L vs E1\_2L})$
    & \textsc{[P]} & 100\% & Headline tokenization gain \\
\addlinespace
\quad (i)~Avoided-tx
    & \textsc{[P]} & \textsc{[P]}\%
    & $\tau_{\text{sell}} + \tau_{\text{buy}}$ burden of E1\_2L moves \\
\qquad of which $\tau_{\text{sell}}$
    & \textsc{[P]} & \textsc{[P]}\%
    & NAR selling commission (6\%) \\
\qquad of which $\tau_{\text{buy}}$
    & \textsc{[P]} & \textsc{[P]}\%
    & Buying closing costs (2.5\%) \\[4pt]
\quad (ii)~Maintained-hedge
    & \textsc{[P]} & \textsc{[P]}\%
    & Continuous-x + pre-buying hedge vs E1\_2L$_\text{NT}$ \\[4pt]
\quad (iii)~Cross-term $\xi$
    & \textsc{[P]} & \textsc{[P]}\%
    & Interaction; expected $\approx 0$ \\
\midrule
v3 Option~3 (no state ext.)
    & $+4.255$\% & (reference) & Previous fire's upper-bound approx \\
\quad of which avoided-tx
    & $+0.816$\% & 19.2\% & From v3 full-grid run \\
\quad of which continuous-x
    & $+3.411$\% & 80.2\% & Liu (2021) territory \\
\quad of which pre-buying hedge
    & $\approx 0$\% & $\approx 0$\% & Absent without state ext. (Option~3 limit) \\
\bottomrule
\end{tabular}
\vspace{4pt}
\begin{minipage}{\linewidth}
\small\textit{Notes:}
Upper panel: v4 Option~1 results (\textsc{[P]} = placeholder).
Lower panel: v3 Option~3 reference numbers (confirmed from server1 run,
commit 186da13).
The maintained-hedge channel in v4 encompasses both the continuous-x
subchannel (present in v3) and the new pre-buying hedge subchannel
enabled by the 6D state; a fourth regime decomposition
($\text{E2\_2L}_{\text{NT}}$) can further separate the two if H1 is confirmed.
\end{minipage}
\end{table}

The maintained-hedge channel in v4 captures two sub-mechanisms:
\begin{enumerate}
    \item \emph{Continuous-$x$ sub-channel} (present in v3): continuous
    fractional ownership of the occupied unit avoids the binary kink at
    $x_\ell = 1$, delivering a Liu~(2021)-style indivisibility relaxation.
    \item \emph{Pre-buying hedge sub-channel} (new in v4): by tracking
    $(x_{A,t-1}, x_{B,t-1})$ in the state, households face a per-period
    buying cost $\tau_{\text{buy}} \cdot \max(\Delta x, 0)$ and can
    pre-accumulate location-$B$ tokens at location $A$, arriving at $B$
    with $x_{B,t-1} > 0$ and avoiding the lump-sum purchase cost.
    The expected per-period premium per unit held is
    $p_{\text{reloc}} \times \tau_{\text{buy}} = 0.06 \times 0.025 = 0.15$\%
    (see Appendix~\ref{app:proof}).
\end{enumerate}

% ─────────────────────────────────────────────────────────────────────────────
\subsection{Falsification Tests}
\label{sec:results:falsification}
% ─────────────────────────────────────────────────────────────────────────────

We conduct three pre-registered falsification tests designed to identify
the pre-buying hedge channel.
These tests were pre-registered in \texttt{docs/welfare\_decomp\_v4.md}
before server1 runs were executed.
Crucially, the analogous tests \emph{failed} under v3 Option~3 (which lacked
the state extension): $\bar{x}_B$ remained zero at all mobility rates,
and the CEV was insensitive to $p_{\text{reloc}}$ and $\rho_{AB}$.

\begin{table}[t]
\centering
\caption{Pre-Registered Falsification Tests
         (Deviations from Baseline)}
\label{tab:falsification}
\small
\begin{tabular}{llcccc}
\toprule
Test & Counterfactual & $\text{CEV}$ & $\bar{x}_B$ at $\ell{=}A$
     & Pass criterion & Result \\
\midrule
(r) Relocation disabled
    & $p_{\text{reloc}} = 0$
    & \textsc{[P]}
    & \textsc{[P]}
    & $\bar{x}_B \to 0$;\; $\Delta\text{CEV} < -0.5$pp
    & \textsc{[P]} \\[4pt]
(m) Perfect correlation
    & $\rho_{AB} = 0.95$
    & \textsc{[P]}
    & \textsc{[P]}
    & $\bar{x}_B$ decreases;\; $\Delta\text{CEV} < 0$
    & \textsc{[P]} \\[4pt]
(q) No buying cost
    & $\tau_{\text{buy}} = 0$
    & \textsc{[P]}
    & \textsc{[P]}
    & Hedge channel $\to 0$;\; $\bar{x}_B \to 0$
    & \textsc{[P]} \\
\bottomrule
\end{tabular}
\vspace{4pt}
\begin{minipage}{\linewidth}
\small\textit{Notes:}
All tests compare to the baseline E2\_2L v4 run.
Pass criterion stated prior to server1 runs.
``Hedge channel'' in test (q) is
$\text{CEV}(\text{E2\_2L}_{\tau=0} \text{ vs E1\_2L}_{\text{NT}}) - \text{CEV}(\text{E2\_2L}_{\text{baseline}} \text{ vs E1\_2L}_{\text{NT}})$.
\textsc{[P]} = placeholder.
\end{minipage}
\end{table}

Test (r) shuts off the relocation shock entirely.
Without future relocation, there is no value to pre-accumulating
location-$B$ tokens while at $A$; the pre-buying hedge channel
should collapse.
If test (r) passes --- i.e., $\bar{x}_B \to 0$ when $p_{\text{reloc}} = 0$
--- it provides strong evidence that the positive $\bar{x}_B$ in the
baseline is driven by the pre-buying motive and not by a model artefact.

Test (m) raises the cross-location correlation to $\rho_{AB} = 0.95$,
nearly eliminating diversification benefit.
When location-$A$ and location-$B$ price processes are nearly identical,
holding $x_B$ at $\ell = A$ provides little hedge against diverging location returns.
The CEV and $\bar{x}_B$ should both decrease relative to the baseline.

Test (q) eliminates the buying cost ($\tau_{\text{buy}} = 0$), which
removes the financial incentive to pre-accumulate $x_B$ (there is no lump-sum
arrival cost to avoid).
The pre-buying hedge channel should collapse to zero.

% ─────────────────────────────────────────────────────────────────────────────
\subsection{Sensitivity Analysis}
\label{sec:results:sensitivity}
% ─────────────────────────────────────────────────────────────────────────────

Figure~\ref{fig:sensitivity_heatmap} depicts the total welfare gain
$\text{CEV}(\text{E2\_2L vs E1\_2L})$ as a function of the two primary
mechanism parameters: the cross-location correlation $\rho_{AB}$ and
the relocation probability $p_{\text{reloc}}$.
The full sweep design follows the pre-registered plan in
\texttt{docs/sensitivity\_grid\_v4.md}; only a subset is shown here.

\begin{figure}[t]
\centering
%% PLACEHOLDER — replace with actual heatmap once sweep results available
\fbox{%
  \begin{minipage}{0.72\linewidth}
  \vspace{2.5cm}
  \centering
  {\sffamily\small [PLACEHOLDER: Figure~2 — CEV sensitivity heatmap.
  X-axis: $\rho_{AB} \in \{0, 0.25, 0.50, 0.75, 0.95\}$.
  Y-axis: $p_{\text{reloc}} \in \{0, 0.02, 0.06, 0.12\}$.
  Heat colour: total CEV(\%), from dark (low) to bright (high).
  Source: sweep outputs from \texttt{scripts/sweep\_rhoAB.sh} and
  \texttt{scripts/sweep\_prelocate.sh} once server1 runs complete.]}
  \vspace{2.5cm}
  \end{minipage}%
}
\caption{Welfare sensitivity to cross-location correlation and relocation rate.
         Each cell reports $\text{CEV}(\text{E2\_2L vs E1\_2L})$ at
         $(\rho_{AB},\, p_{\text{reloc}})$.
         The mechanism prediction is a positive gradient in $p_{\text{reloc}}$
         (higher mobility $\Rightarrow$ higher hedge value) and a negative
         gradient in $\rho_{AB}$ (more correlated locations $\Rightarrow$
         smaller diversification benefit).}
\label{fig:sensitivity_heatmap}
\end{figure}

Table~\ref{tab:sensitivity} reports selected cross-sections of the
sensitivity grid for other key parameters.

\begin{table}[t]
\centering
\caption{Sensitivity of Headline CEV to Key Parameters
         (Total CEV, \%; baseline value on diagonal)}
\label{tab:sensitivity}
\small
\begin{tabular}{lcccccc}
\toprule
Parameter & & \multicolumn{5}{c}{Values} \\
\midrule
Risk aversion $\gamma$
    & & 3 & \textbf{5} & 8 & & \\
\quad $\text{CEV}(\text{E2\_2L vs E1\_2L})$
    & & \textsc{[P]} & \textsc{[P]} & \textsc{[P]} & & \\
\addlinespace
Buying cost $\tau_{\text{buy}}$
    & & 0\% & 1\% & \textbf{2.5\%} & 4\% & 6\% \\
\quad $\text{CEV}(\text{E2\_2L vs E1\_2L})$
    & & \textsc{[P]} & \textsc{[P]} & \textsc{[P]} & \textsc{[P]} & \textsc{[P]} \\
\addlinespace
Relocation rate $p_{\text{reloc,work}}$
    & & 0 & 0.02 & \textbf{0.06} & 0.12 & \\
\quad $\text{CEV}(\text{E2\_2L vs E1\_2L})$
    & & \textsc{[P]} & \textsc{[P]} & \textsc{[P]} & \textsc{[P]} & \\
\quad $\bar{x}_B$ at $\ell{=}A$
    & & \textsc{[P]} & \textsc{[P]} & \textsc{[P]} & \textsc{[P]} & \\
\addlinespace
Cross-location corr.\ $\rho_{AB}$
    & & 0 & 0.25 & \textbf{0.50} & 0.75 & 0.95 \\
\quad $\text{CEV}(\text{E2\_2L vs E1\_2L})$
    & & \textsc{[P]} & \textsc{[P]} & \textsc{[P]} & \textsc{[P]} & \textsc{[P]} \\
\addlinespace
\multicolumn{7}{l}{\textit{Asymmetric robustness (solver v4 extension, fire 22)}} \\
Location-$B$ return premium $\Delta\mu_B$
    & & $-1\%$ & \textbf{0\%} & $+1\%$ & & \\
\quad $\text{CEV}(\text{E2\_2L vs E1\_2L})$
    & & \textsc{[P]} & \textsc{[P]} & \textsc{[P]} & & \\
\quad $\bar{x}_B$ at $\ell{=}A$
    & & \textsc{[P]} & \textsc{[P]} & \textsc{[P]} & & \\[2pt]
Directional mobility ($p_{A\to B}$, $p_{B\to A}$)
    & & (0.06, 0.06) & (0.10, 0.03) & (0.10, 0.10) & & \\
\quad $\text{CEV}(\text{E2\_2L vs E1\_2L})$
    & & \textsc{[P]} & \textsc{[P]} & \textsc{[P]} & & \\
\addlinespace
\multicolumn{7}{l}{\textit{Mortgage activation (LTV constraint)}} \\
LTV ceiling $\theta_{\max}$
    & & \textbf{0} & 0.50 & 0.80 & & \\
\quad $\text{CEV}(\text{E2\_2L vs E1\_2L})$
    & & \textsc{[P]} & \textsc{[P]} & \textsc{[P]} & & \\
\bottomrule
\end{tabular}
\vspace{4pt}
\begin{minipage}{\linewidth}
\small\textit{Notes:}
Bold entries mark the baseline value.
One parameter varies at a time; all others held at baseline.
Sweep scripts: \texttt{scripts/sweep\_rhoAB.sh},
\texttt{scripts/sweep\_prelocate.sh}, \texttt{scripts/sweep\_txcost.sh},
\texttt{scripts/sweep\_asymmetric.sh} (location-$B$ return premium and directional mobility),
\texttt{scripts/sweep\_mortgage.sh} (LTV).
$\Delta\mu_B = \mu_{h,B} - \mu_{h,A}$: the excess Jensen-corrected log-mean return of
location $B$ relative to $A$.
Positive $\Delta\mu_B$ strengthens the pre-buying motive for $x_B$ by adding a return
premium to the tau\_buy-saving incentive.
\textsc{[P]} = placeholder pending server1 sweep runs.
\end{minipage}
\end{table}

Two monotonicity predictions follow from the model mechanism.
First, for $\tau_{\text{buy}}$: as the buying cost rises, the pre-accumulation
motive strengthens and $\bar{x}_B$ should increase, lifting the hedge sub-channel.
Second, for $p_{\text{reloc}}$: higher mobility raises the expected value of
pre-accumulated $x_B$ holdings; at $p_{\text{reloc}} = 0$ the hedge channel
collapses (falsification test r).
If these monotonicity patterns hold in the sweep results, they constitute
additional evidence for the pre-buying hedge mechanism.

% ─────────────────────────────────────────────────────────────────────────────
\subsection{Summary of Mechanism Evidence}
\label{sec:results:summary}
% ─────────────────────────────────────────────────────────────────────────────

The results establish \textsc{[P]} as the headline lifetime welfare gain from
tokenization under the baseline calibration.
The gain decomposes additively into an avoided-transaction-cost channel
(\textsc{[P]}\%) and a maintained-hedge channel (\textsc{[P]}\%), with
a near-zero cross-term.
The maintained-hedge channel is itself divisible into the continuous-$x$
sub-channel present in prior work (Liu, 2021) and the novel pre-buying hedge
sub-channel unlocked by the 6D state extension.

The three falsification tests provide mechanism identification.
Under all three counterfactuals --- zero relocation, high cross-location
correlation, zero buying cost --- the pre-buying hedge channel should
attenuate or collapse.
Confirmation of these patterns will establish that the welfare gain reported
in Table~\ref{tab:cev_baseline} reflects the model's structural mechanism
rather than a calibration artefact.
\textsc{[Fill in pass/fail outcomes from output/diagnostics/p6\_option1\_decomposition.md.]}
